%% ============================================================
%% CodeChange20260520.tex — ARES 코드 수정 보고서 (합본)
%% 2026-05-19 ~ 2026-05-20
%% BLE 파이프라인 최적화, 부팅 안정화, DC 모터 대칭 출력, HTTPS 사유
%% XeLaTeX 컴파일: xelatex -shell-escape CodeChange20260520.tex
%% ============================================================
\documentclass[11pt, oneside]{article}

%% ----- 한국어 / 폰트 -----
\usepackage{kotex}
\usepackage{fontspec}
\defaultfontfeatures{Mapping=}
%% 사용 환경에 맞게 수정. macOS: /Users/<USER>/Library/Fonts/
\newcommand{\userfontpath}{C:/Users/young/AppData/Local/Microsoft/Windows/Fonts/}

\setmainfont{KoPub Batang}[
  Path=\userfontpath, Scale=0.92, Mapping=, LetterSpace=-2.0,
  UprightFont = KoPub Batang Medium.ttf,
  BoldFont    = KoPub Batang Bold.ttf
]
\setsansfont{KoPub Dotum}[
  Path=\userfontpath, Scale=0.92, Mapping=, LetterSpace=-2.0,
  UprightFont = KoPub Dotum Medium.ttf,
  BoldFont    = KoPub Dotum Bold.ttf
]
\setmonofont{D2Coding}[
  Path=\userfontpath, Scale=0.88,
  Extension = .ttc,
  UprightFont = D2Coding-Ver1.3.2-20180524-all
]
\setmainhangulfont{KoPub Batang}[
  Path=\userfontpath, Scale=0.92, Mapping=, LetterSpace=-2.0,
  UprightFont = KoPub Batang Medium.ttf,
  BoldFont    = KoPub Batang Bold.ttf
]
\setsanshangulfont{KoPub Dotum}[
  Path=\userfontpath, Scale=0.92, Mapping=, LetterSpace=-2.0,
  UprightFont = KoPub Dotum Medium.ttf,
  BoldFont    = KoPub Dotum Bold.ttf
]
\setmonohangulfont{KoPub Batang}[
  Path=\userfontpath, Scale=0.92, LetterSpace=-2.0,
  UprightFont = KoPub Batang Medium.ttf,
  BoldFont    = KoPub Batang Bold.ttf
]

%% ----- 레이아웃 -----
\usepackage{geometry}
\geometry{
  paperwidth=190mm, paperheight=240mm,
  top=24mm, bottom=28mm, inner=24mm, outer=22mm,
  headheight=14pt, headsep=8mm, footskip=28pt
}
\linespread{1.12}
\setlength{\parskip}{0.85em}
\setlength{\parindent}{0pt}
\raggedbottom
\setcounter{tocdepth}{2}
\setcounter{secnumdepth}{2}

%% ----- 색상 / 박스 / 코드 -----
\usepackage{xcolor}
\definecolor{deepblue}{RGB}{22, 56, 100}
\definecolor{accentcopper}{RGB}{200, 105, 50}
\definecolor{labcolor}{RGB}{40, 140, 125}
\definecolor{cautioncolor}{RGB}{195, 85, 70}
\definecolor{rulegray}{RGB}{210, 214, 220}
\definecolor{codebackground}{RGB}{249, 250, 252}

\usepackage{minted}
\usemintedstyle{friendly}
\setminted{
  fontsize=\footnotesize,
  bgcolor=codebackground,
  frame=leftline,
  framerule=0.8pt,
  rulecolor=\color{accentcopper},
  breaklines=true,
  breakanywhere=true,
  tabsize=2,
  linenos=false,
  autogobble=true,
}

\usepackage{tcolorbox}
\tcbuselibrary{skins, breakable}

\newtcolorbox{keypoint}{
  enhanced, breakable, sharp corners, arc=0pt, frame hidden, boxrule=0pt,
  colback=accentcopper!6,
  borderline west={2pt}{0pt}{accentcopper},
  top=6pt, bottom=6pt, left=12pt, right=10pt,
  before skip=1em, after skip=0.9em,
}
\newtcolorbox{caution}{
  enhanced, breakable, sharp corners, arc=0pt, frame hidden, boxrule=0pt,
  colback=cautioncolor!6,
  borderline west={2pt}{0pt}{cautioncolor},
  top=6pt, bottom=6pt, left=12pt, right=10pt,
  before skip=1em, after skip=0.9em,
}

%% ----- 표 -----
\usepackage{booktabs}
\usepackage{tabularx}
\usepackage{array}

%% ----- 섹션 스타일 (간결) -----
\usepackage{titlesec}
\titleformat{\section}[block]
  {\sffamily\bfseries\Large\color{deepblue}}{\thesection}{0.6em}{}
  [\vspace{2pt}{\color{rulegray}\hrule height 0.5pt}\vspace{4pt}]
\titleformat{\subsection}
  {\sffamily\bfseries\normalsize\color{accentcopper}}{\thesubsection}{0.6em}{}

%% ----- 헤더 / 푸터 -----
\usepackage{fancyhdr}
\pagestyle{fancy}
\fancyhf{}
\renewcommand{\headrulewidth}{0pt}
\renewcommand{\footrulewidth}{0pt}
\fancyfoot[L]{\sffamily\small\color{deepblue}\bfseries ARES BLE 파이프라인 · 부팅 안정화 · DC 모터 · HTTPS 사유}
\fancyfoot[R]{\sffamily\small\thepage}

\usepackage{hyperref}
\hypersetup{
  colorlinks=true, linkcolor=deepblue, urlcolor=deepblue,
  pdftitle={ARES 코드 수정 보고서 — 2026-05-19~2026-05-20 합본},
  pdfauthor={(주)코리아사이언스, 동명대학교 게임학부},
}

\renewcommand{\contentsname}{목차}
\renewcommand{\figurename}{그림}
\renewcommand{\tablename}{표}

\title{%
  {\sffamily\bfseries\fontsize{20}{24}\selectfont\color{deepblue}ARES 코드 수정 보고서}\\[4pt]
  {\rmfamily\normalsize\color{accentcopper}BLE 명령 파이프라인 최적화 · 부팅 안정화 · DC 모터 대칭 출력 · HTTPS 사용 사유}\\[2pt]
  {\rmfamily\footnotesize\color{rulegray}\textit{2026-05-19 \textendash{} 2026-05-20 합본}}
}
\author{\sffamily (주)코리아사이언스\quad/\quad 동명대학교 게임학부}
\date{2026년 5월 19일 \textendash{} 5월 20일}

\begin{document}
\maketitle
\thispagestyle{plain}

\section{개요}

본 보고서는 2026년 5월 19일과 5월 20일에 적용된 ARES 펌웨어 및 웹 코드의
변경 사항을 하나의 합본으로 정리한 자료이다. 두 일자의 변경은 형식적으로는
구분되지만 실질적으로 연결되어 있으며, 5월 19일에 적용한 BLE 명령
파이프라인 최적화의 \emph{재적용 과정}에서 부팅 단계 크래시와 SYS\_SET
핸들러 누락이 드러났고, 같은 흐름에서 DC 모터 출력 비대칭과 Web Bluetooth
호스팅 환경 의존성도 정리되었다.

전체 변경의 큰 줄기는 다음과 같다.

\begin{enumerate}
  \item \textbf{(2026-05-19)} BLE 명령 파이프라인 최적화. ``LED 1번 켜기
        $\rightarrow$ LED 2번 켜기''와 같이 연속된 출력 명령 사이의
        체감 지연(약 250\,ms)을 명령 부류별 응답 정책 분리로 약 40\,ms
        수준까지 단축.
  \item \textbf{(2026-05-20)} 같은 최적화를 재적용하면서 발견된 \emph{펌웨어
        부팅 단계의 OLED 크래시}와 \emph{SYS\_SET 명령 핸들러 누락}을
        교정. 또한 DC 모터의 \texttt{DIR+PWM} 단일 변조 토폴로지에서
        오는 전·후진 출력 비대칭을 듀얼 PWM brake-modulated 구성으로
        근원에서 해소.
  \item \textbf{(2026-05-20)} 학습용 웹앱 \texttt{Web/main.html}이
        \texttt{file://}에서 Web Bluetooth가 동작하지 않는 원인 식별.
        보안 컨텍스트의 문제가 아니라 \emph{ES module의 같은-출처
        정책}이 원인이며, 향후 학생 배포 시 HTTPS 호스팅(또는
        \texttt{localhost})이 사실상 강제됨을 명시.
\end{enumerate}

\begin{keypoint}
2026-05-19의 작업이 \emph{기능적 개선}이라면, 2026-05-20의 작업은 그
개선이 \emph{재현 가능하게 동작}하도록 만든 안정화이다. 같은 코드라도
하나의 회피 가능한 크래시 때문에 정상이 비정상으로 보이고, 같은 모터라도
PWM 와이어링 한 가지가 사용자 경험을 크게 좌우한다.
\end{keypoint}

\section{2026-05-19 --- BLE 명령 파이프라인 최적화}

ARES 로버의 블록 코딩 사용 중 ``LED 1번 켜기 $\rightarrow$ LED 2번 켜기''와 같이
연속된 명령 사이에 체감 가능한 지연이 발생하던 문제를 분석하여, 명령당
누적 약 250\,ms에 달하던 지연을 출력 명령 기준 약 40\,ms 수준으로 줄였다.
변경 대상은 \texttt{Web/commandexecutor.js}와 \texttt{Pico/main.py}이며,
명령의 의미에 따라 응답 라운드트립을 \emph{선택적으로} 생략하는 정책을
양쪽에 일관되게 도입하였다.

\subsection{문제 진단}

단일 LED 명령 한 번이 종단 간에 소요하는 시간 예산은 다음과 같이
구성되어 있었다.

\begin{table}[h]
\centering
\small
\begin{tabularx}{\linewidth}{@{}llXr@{}}
\toprule
\# & 위치 & 원인 & 시간 \\
\midrule
1 & \texttt{commandexecutor.js:199} & \texttt{sendData(cmd, true)}로 모든 명령에 응답 대기 & 라운드트립 \\
2 & \texttt{bluetooth.js:396} & \texttt{CHUNK\_DELAY=50ms} (\texttt{constants.js:14}) & 50ms \\
3 & \texttt{main.py:86} & \texttt{RECEIVE\_TIMEOUT\_MS=50ms} 무조건 대기 & 50ms \\
4 & \texttt{commandexecutor.js:213} & 모든 명령 후 \texttt{setTimeout(100)} & 100ms \\
5 & BLE GATT  & connection interval (HM-10 기본) & 30--70ms \\
\bottomrule
\end{tabularx}
\caption{단일 LED 명령의 누적 지연(개선 전) --- 합계 약 250\,ms.}
\end{table}

\begin{keypoint}
정상 동작이라도 누적 지연이 \textbf{약 250\,ms / 명령}에 달했다. 카운트
다운, 파티 조명처럼 짧은 간격의 LED 시퀀스가 ``끊기는'' 체감의 주된
원인이다.
\end{keypoint}

\subsection{변경 정책 --- 명령별 분리}

명령을 두 부류로 나누어 응답 대기 여부와 cooldown을 다르게 적용한다.

\begin{table}[h]
\centering
\small
\begin{tabularx}{\linewidth}{@{}lXl@{}}
\toprule
부류 & 명령 & 응답 대기 \\
\midrule
Fire-and-forget & \texttt{LED\_ON}, \texttt{LED\_OFF}, \texttt{MAIN\_LED\_*}, \texttt{[v0 v1 v2 v3 v4]} (LED 패턴),
\texttt{MSG}, \texttt{CLEAR\_DISPLAY}, \texttt{SERVO\_FORWARD}/\texttt{BACKWARD}/\texttt{LEFT}/\texttt{RIGHT}/\texttt{STOP} (연속),
\texttt{DC\_FORWARD}/\texttt{BACKWARD}/\texttt{STOP} (연속), \texttt{GUN\_FIRE} & 아니오 \\
\midrule
응답 대기 유지 & \texttt{BUZZER\_ON,F,D} (D초 blocking), \texttt{SERVO\_t*}/\texttt{DC\_t*} (N초 blocking),
\texttt{SLEEP,N} (N초 blocking), \texttt{DISTANCE}, \texttt{MAGNET}, \texttt{PING} & 예 \\
\bottomrule
\end{tabularx}
\caption{명령 부류와 응답 대기 정책. Pico에서 blocking으로 처리되거나
값을 반환하는 명령은 응답 대기를 유지하여 타이밍과 값 정합성을 보존한다.}
\end{table}

\subsection{Web 측 변경 --- \texttt{Web/commandexecutor.js}}

\subsubsection*{핵심 코드}

\begin{minted}{javascript}
// 즉시 완료 명령 — Pico가 blocking 처리하지 않으므로 응답 대기 불필요.
// 시간 의존적(BUZZER_ON, SERVO_t*, DC_t*, SLEEP) 또는 값 반환(DISTANCE,
// MAGNET, PING)은 이 집합에 넣지 말 것. 응답 대기로 두어야 타이밍/값이
// 보장된다.
FIRE_AND_FORGET_HEADS: new Set([
  'LED_ON', 'LED_OFF',
  'MAIN_LED_ON', 'MAIN_LED_OFF',
  'MSG', 'CLEAR_DISPLAY',
  'SERVO_FORWARD', 'SERVO_BACKWARD', 'SERVO_LEFT', 'SERVO_RIGHT', 'SERVO_STOP',
  'DC_FORWARD', 'DC_BACKWARD', 'DC_STOP',
  'GUN_FIRE',
]),

_isFireAndForget(command) {
  if (command.startsWith('[')) return true;     // LED 패턴 [v0 v1 v2 v3 v4]
  const head = command.split(',')[0];
  return this.FIRE_AND_FORGET_HEADS.has(head);
},

async sendCommand(command) {
  if (!state.isExecuting) {
    Logger.add('[중단] 실행이 중단되었습니다', 'warning');
    return;
  }
  BluetoothManager.updateStatus('명령 실행 중...', STATUS_COLORS.ORANGE);

  const fireAndForget = this._isFireAndForget(command);

  try {
    await BluetoothManager.sendData(command, !fireAndForget);
    if (DEBUG) Logger.add(`[완료] ${command}`, 'info');
  } catch (error) {
    // ... 에러 처리 ...
  }

  // 송신 간격 가드.
  // - fire-and-forget: UART(9600 baud)/BLE 버퍼 오버플로우 방지 최소 마진.
  // - 응답 대기: 이미 라운드트립을 거쳤으므로 가드 단축.
  const cooldown = fireAndForget ? 40 : 20;
  await new Promise(resolve => setTimeout(resolve, cooldown));
}
\end{minted}

\subsubsection*{효과}

\begin{itemize}
  \item 출력 명령(LED, SERVO 연속, DC 연속 등): BLE 라운드트립과 100\,ms
        cooldown이 모두 사라짐 $\rightarrow$ 명령당 약 40\,ms.
  \item 응답 대기 명령: cooldown만 100\,ms$\rightarrow$20\,ms 단축. 의미적
        동기화는 그대로 유지.
\end{itemize}

\subsection{Pico 측 변경 --- \texttt{Pico/main.py}}

웹의 분류와 \emph{동일 기준}으로 Pico에서도 응답 송신 여부를 결정한다.
또한 단일 chunk 명령에 대한 무조건 50\,ms 대기를 제거한다.

\subsubsection*{\texttt{NO\_RESPONSE\_CMDS} 및 \texttt{\_needs\_response}}

\begin{minted}{python}
class AresRover:
    # 응답 송신 생략 명령 (Web/commandexecutor.js의 FIRE_AND_FORGET_HEADS와 동기화).
    # 즉시 완료되는 출력 명령은 응답 라운드트립을 생략하여 처리량을 높이고,
    # 응답 매칭 어긋남(LED_ON 응답이 다음 명령 슬롯에 잘못 들어가는 현상)을 차단한다.
    NO_RESPONSE_CMDS = (
        "LED_ON", "LED_OFF", "MAIN_LED_ON", "MAIN_LED_OFF",
        "MSG", "CLEAR_DISPLAY",
        "SERVO_FORWARD", "SERVO_BACKWARD", "SERVO_LEFT", "SERVO_RIGHT", "SERVO_STOP",
        "DC_FORWARD", "DC_BACKWARD", "DC_STOP",
        "GUN_FIRE",
    )

    def _needs_response(self, data):
        """응답 송신 여부. fire-and-forget 명령은 False (NO_RESPONSE_CMDS 참조)."""
        if data.startswith("["):           # LED 패턴 [v0 v1 v2 v3 v4]
            return False
        head = data.split(",", 1)[0]
        return head not in self.NO_RESPONSE_CMDS
\end{minted}

\subsubsection*{\texttt{\_read\_uart\_line()} 폴링 방식 교체}

기존 구현은 매 호출마다 \texttt{utime.sleep\_ms(RECEIVE\_TIMEOUT\_MS)}로
50\,ms를 무조건 대기한 뒤 두 번째 read를 수행했다. newline이 이미
도착한 단일 chunk 명령에도 적용되어 처리량 천장이 약 16--17 cmd/s에
묶여 있었다. 폴링으로 교체하여 newline 발견 즉시 종료하도록 한다.

\begin{minted}{python}
def _read_uart_line(self):
    """UART에서 완전한 라인 읽기.
    newline 발견 즉시 종료 — 단일 chunk 명령에서 무조건 50ms를 기다리지 않는다.
    멀티 chunk 명령(LED 패턴, SYS_SET 등)은 RECEIVE_TIMEOUT_MS까지 폴링하여 안전.
    """
    if not self.uart.any() and not self.rx_buffer:
        return None

    start = utime.ticks_ms()
    while utime.ticks_diff(utime.ticks_ms(), start) < RECEIVE_TIMEOUT_MS:
        if self.uart.any():
            chunk = self.uart.read()
            if chunk:
                self.rx_buffer += chunk.decode('utf-8', 'ignore')

            # 버퍼 오버플로우 방지
            if len(self.rx_buffer) > MAX_BUFFER_SIZE:
                print("[UART] 버퍼 오버플로우, 초기화")
                self.rx_buffer = ""
                return None

            # newline 즉시 종료
            if '\n' in self.rx_buffer:
                newline_idx = self.rx_buffer.index('\n')
                complete_line = self.rx_buffer[:newline_idx].strip()
                self.rx_buffer = self.rx_buffer[newline_idx + 1:]
                return complete_line
        else:
            utime.sleep_ms(2)

    # 타임아웃: 라인 종료자 없는 부분 데이터 반환
    if self.rx_buffer:
        data = self.rx_buffer.strip()
        if data and not data.endswith(','):
            self.rx_buffer = ""
            return data

    return None
\end{minted}

\subsubsection*{\texttt{run()}의 응답 조건부 송신}

\begin{minted}{python}
# 변경 전
else:
    response = self._process_command(data)
    self._send_response(response)

# 변경 후
else:
    response = self._process_command(data)
    if self._needs_response(data):
        self._send_response(response)
\end{minted}

\subsection{기대 효과 (정량)}

\begin{table}[h]
\centering
\small
\begin{tabular}{@{}lcc@{}}
\toprule
명령 종류 & 변경 전 & 변경 후 \\
\midrule
\texttt{LED\_ON} 등 fire-and-forget & 약 250\,ms & 약 40\,ms \\
\texttt{BUZZER\_ON,F,D} 등 blocking & $D + 250$\,ms & $D + 20$\,ms \\
\texttt{DISTANCE}/\texttt{MAGNET} 등 값 반환 & 약 250\,ms & RTT $+$ 20\,ms \\
Pico 처리량(단일 chunk) & $\approx 16$--$17$ cmd/s & $\approx 100^+$ cmd/s \\
\bottomrule
\end{tabular}
\caption{명령 부류별 종단 간 지연 비교 (단위: ms).}
\end{table}

추가로 응답 매칭 어긋남(이전 fire-and-forget 명령의 응답이 다음
응답 대기 명령의 \texttt{pendingResolve} 슬롯에 잘못 들어가는 코너
케이스) 위험이 \emph{구조적으로} 제거된다. Pico에서 응답 자체를
보내지 않으므로 침범할 응답이 존재하지 않는다.

\begin{caution}
변경된 두 파일은 \textbf{같이 갱신}되어야 한다. 한쪽만 갱신되면:
\begin{itemize}
  \item Web만 갱신: 출력 명령에 응답 대기를 안 하지만 Pico는 여전히
        응답을 보냄 $\rightarrow$ BLE notify 누적, 응답 매칭 어긋남 위험 잔존.
  \item Pico만 갱신: Web이 출력 명령에 응답 대기를 하지만 Pico가 응답을
        안 보냄 $\rightarrow$ 매 명령 5초 timeout.
\end{itemize}
\end{caution}

\section{2026-05-20 --- 재적용과 부팅 안정화}

\subsection{재적용 배경}

전일 적용된 BLE 파이프라인 최적화(\texttt{7e5b93f})는 적용 직후 BLE
연결 자체는 이루어지지만 \texttt{GET\_STATE} 등 모든 명령이 응답하지
않는 현상이 보고되어 \texttt{110bb27}로 revert 되어 있었다.

재현 후 원인 분석한 결과, BLE 파이프라인 최적화 자체는 문제가 없었다.
진짜 원인은 \texttt{Pico/main.py}의 \texttt{boot()}에서 OLED 초기화
실패 시 \texttt{robot.oled = None} 상태인데 \texttt{None.fill(0)}이
호출되어 \texttt{AttributeError}가 발생하고, 이 예외가 메인 루프 진입을
원천적으로 막아버린 데 있었다. 즉 \emph{펌웨어 부팅 단계의 무관한
크래시}가 BLE 파이프라인 변경의 효과를 가린 것이다.

\subsection{관찰의 단서 --- BLE 연결은 되는데 응답이 없다}

BLE 모듈(BT05/HM-10)은 Pico의 UART에 연결된 외부 모듈로, BLE
페어링/연결은 모듈 자체 펌웨어가 처리하고 Pico는 UART로 데이터만
주고받는다. 따라서 Pico의 \texttt{main.py}가 boot에서 죽어도 BLE 링크는
멀쩡하다. 외부에서 보이는 증상은 ``연결은 되는데 어떤 명령도
응답하지 않는다''로, 이는 정확히 \emph{Pico가 메인 루프에 진입조차
못한} 상태의 결과이다.

\subsection{수정 --- OLED None 가드}

\begin{minted}{python}
# 변경 전 (Pico/main.py boot())
self.robot.oled.fill(0)
self.robot.oled.text("ARES READY", 0, 0)
self.robot.oled.show()

# 변경 후
if self.robot.oled is not None:
    self.robot.oled.fill(0)
    self.robot.oled.text("ARES READY", 0, 0)
    self.robot.oled.show()
\end{minted}

가드 추가 후 OLED 모듈이 비활성(또는 I2C 초기화 실패) 상태여도
\texttt{boot()}가 완료되며, 메인 루프가 정상 진입한다. 이 가드는 다른
하드웨어 모듈에 적용된 \texttt{if robot.module:} 패턴과 일관된
방어 코드이다.

가드 적용 후 5월 19일의 BLE 최적화를 동일 형태로 재적용하였고
(\texttt{600243a}), 정상 동작을 확인하였다.

\section{2026-05-20 --- SYS\_SET 핸들러 누락 수정}

\subsection{증상}

대시보드에서 ``최대 속도'' 슬라이더를 낮추고 ``Pico에 업로드'' 버튼을
누르면 ``설정이 Pico에 업로드되었습니다!'' 라는 알림이 떴지만, 다른
탭으로 갔다가 돌아오면 슬라이더가 다시 80\%로 복귀했다.

\subsection{원인}

\texttt{Pico/process\_data.py}의 명령 디스패처에서:

\begin{minted}{python}
elif data.startswith("SYS_SET"):
    return self._handle_sys_set(data)
\end{minted}

호출 대상인 \texttt{\_handle\_sys\_set} 메서드가 \emph{정의되어 있지
않았다}. 메서드 헤더(\texttt{def \_handle\_sys\_set(self, data):})가
빠져 있어, 핸들러 본문이 직전의 \texttt{\_handle\_set\_module}의
\texttt{return} 뒤에 떨어져 \emph{도달 불가능한 코드(dead code)} 가
되어 있었다. SYS\_SET 명령이 도달하면 \texttt{AttributeError}가 발생
하여 \texttt{main.py}의 예외 처리에 잡혀 \texttt{"ERROR"}가 반환되고,
\texttt{system.json}은 갱신되지 않았다. Web 측은 응답을 확인하지 않고
즉시 alert를 띄우므로 사용자에게는 ``저장된 것처럼'' 보였다.

\subsection{수정}

\begin{minted}{python}
def _handle_sys_set(self, data):
    """시스템 설정 저장: SYS_SET,max_speed,col_dist,auto_stop,name"""
    try:
        parts = data.split(',')
        if len(parts) >= 5:
            max_speed = parts[1]
            col_dist  = parts[2]
            auto_stop = parts[3]
            name      = ",".join(parts[4:])
            sys_config.save_config(max_speed, col_dist, auto_stop, name)
            return 1
    except Exception as e:
        print(f"SYS_SET 오류: {e}")
    return 0
\end{minted}

수정 후 슬라이더 값이 \texttt{system.json}에 영구 저장되고, Pico 재부팅
후에도 유지된다.

\begin{caution}
Web 측 \texttt{saveSystemSettings()}는 SYS\_SET 응답을 검증하지 않고
무조건 성공 alert를 띄운다. 후속 개선 항목으로 \texttt{sendData(cmd,
true)}로 응답 대기 후 결과를 분기하는 작업을 분리하여 진행한다.
\end{caution}

\section{2026-05-20 --- DC 모터 H-브릿지 비대칭 진단과 듀얼 PWM 대칭화}

\subsection{증상의 변천}

\begin{enumerate}
  \item 초기: 전진이 후진보다 명백히 약함(특히 \texttt{max\_speed}가
        높을수록 격차가 큼).
  \item 단순 PWM 반전 적용 후: 80\,\%에서는 양방향 정상이나, 슬라이더
        55\,\%\,이하에서 후진이 시동되지 않음.
  \item 후진에 시동 데드존 lift 적용(56\,\%): 시동은 되지만 후진 토크가
        56--60\,\% 구간에서 약함.
  \item 모터 리드를 물리적으로 뒤집는 진단 실험: 회전 방향만 반대로
        바뀌고 전·후진 속도 특성은 함수 호출 기준 그대로 유지됨.
\end{enumerate}

\subsection{진단 결과 --- 비대칭은 모터가 아니라 드라이버 입력 토폴로지}

기존 코드는 \texttt{DIR + PWM} 한 쌍으로 H-브릿지를 구동하였다. 이때
드라이버 칩의 \texttt{IN1}/\texttt{IN2} 두 입력 중 한쪽은 디지털 DIR로
고정되고 다른 쪽만 PWM이 들어가므로, 두 회전 방향이 \emph{서로 다른
PWM 변조 모드}로 동작한다.

\begin{table}[h]
\centering
\small
\begin{tabularx}{\linewidth}{@{}lXX@{}}
\toprule
모드 & 동작 & 특성 \\
\midrule
\texttt{DIR=1}, IN2 PWM (전진) & drive$\leftrightarrow$\textbf{brake} 변조 & 권선 전류 연속, 시동 우수, 듀티가 작을수록 토크 큼 \\
\texttt{DIR=0}, IN2 PWM (후진) & drive$\leftrightarrow$\textbf{coast} 변조 & 권선 전류 불연속, 시동 데드존 존재, 약 56\,\% 이하에서 시동 불가 \\
\bottomrule
\end{tabularx}
\caption{단일 \texttt{DIR+PWM} 토폴로지에서 회전 방향별 PWM 변조 모드.}
\end{table}

모터 리드를 뒤집어도 함수 단위의 비대칭이 그대로 유지된 실험 결과가
이 가설을 결정적으로 확증한다. 비대칭은 \emph{드라이버 입력 측 패턴}에
귀속되며, 모터 권선이나 기어박스의 마찰 비대칭은 아니다.

\subsection{듀얼 PWM brake-modulated 구동}

RP2040에서 \texttt{GP8}/\texttt{GP9}는 같은 PWM 슬라이스(PWM4)의 A/B
채널이다. 두 핀을 모두 PWM으로 설정한 뒤, 한쪽을 \texttt{PWM\_MAX}로
고정하고 반대쪽을 \texttt{PWM\_MAX - speed}로 변조하면 양방향 모두
brake-modulated가 되어 토크 곡선이 대칭화된다.

\begin{minted}{python}
# Pico/dcmotor.py
from machine import Pin, PWM
from pins import DCMOTOR_DIR_PIN, DCMOTOR_PWM_PIN

IN1_PIN = DCMOTOR_DIR_PIN  # GP8 (PWM4A)  — 이름은 레거시 호환 유지
IN2_PIN = DCMOTOR_PWM_PIN  # GP9 (PWM4B)

PWM_FREQUENCY = 20000
PWM_STOP = 0
PWM_MAX  = 65535


class KSDCMotor:
    """DC 모터 컨트롤러 (양방향 brake-modulated H-브릿지)."""

    def __init__(self):
        self.IN1 = PWM(Pin(IN1_PIN))
        self.IN1.freq(PWM_FREQUENCY)
        self.IN1.duty_u16(PWM_STOP)

        self.IN2 = PWM(Pin(IN2_PIN))
        self.IN2.freq(PWM_FREQUENCY)
        self.IN2.duty_u16(PWM_STOP)

        self._is_running = False

    def dc_forward(self, speed):
        """전진: IN1 고정 HIGH, IN2를 (PWM_MAX-speed)로 변조."""
        if speed <= 0:
            self.dc_stop(); return
        self.IN1.duty_u16(PWM_MAX)
        self.IN2.duty_u16(PWM_MAX - speed)
        self._is_running = True

    def dc_backward(self, speed):
        """후진: IN2 고정 HIGH, IN1을 (PWM_MAX-speed)로 변조."""
        if speed <= 0:
            self.dc_stop(); return
        self.IN1.duty_u16(PWM_MAX - speed)
        self.IN2.duty_u16(PWM_MAX)
        self._is_running = True

    def dc_stop(self):
        """정지 (coast: 둘 다 LOW)."""
        self.IN1.duty_u16(PWM_STOP)
        self.IN2.duty_u16(PWM_STOP)
        self._is_running = False
\end{minted}

\subsection{기대 효과}

\begin{table}[h]
\centering
\small
\begin{tabularx}{\linewidth}{@{}lXXX@{}}
\toprule
슬라이더 & 전진 실효 듀티 & 후진 실효 듀티 & 비고 \\
\midrule
100\,\% & 100\,\% & 100\,\% & 풀파워, 양방향 동일 \\
80\,\%  & 80\,\%  & 80\,\%  & 동일 \\
50\,\%  & 50\,\%  & 50\,\%  & 동일 \\
20\,\%  & 20\,\%  & 20\,\%  & 동일 (시동은 모터 마찰 한계에 의존) \\
0\,\%   & 0       & 0       & coast \\
\bottomrule
\end{tabularx}
\caption{듀얼 PWM brake-modulated 구동의 양방향 대칭 출력.}
\end{table}

\begin{keypoint}
시동 데드존(\texttt{BACKWARD\_DEADZONE}) lift 같은 보정 hack이 더는
필요하지 않다. \emph{소프트웨어로 우회}하던 비대칭을 \emph{펌웨어
구성 변경}으로 \emph{근원에서} 제거한 사례이다. 회로 배선이나 드라이버
교체 없이 펌웨어만으로 해결되었다는 점이 특히 가치 있다.
\end{keypoint}

\subsection{잔여 사항}

\begin{itemize}
  \item \texttt{pins.py} 및 \texttt{system\_config.py}의 핀 이름
        (\texttt{DCMOTOR\_DIR\_PIN}, \texttt{pin\_dcmotor\_dir})은
        의미상 ``DIR''이 아니지만 외부 API(\texttt{SET\_PIN} 명령
        등)와의 호환을 위해 레거시 명칭을 유지했다. 향후 정리 시
        \texttt{IN1}/\texttt{IN2} 또는 \texttt{A}/\texttt{B}로
        명명 통일이 권장된다.
  \item 모터의 정적 마찰 한계로 약 20\,\% 이하의 매우 낮은 듀티에서는
        여전히 회전하지 않을 수 있다. 이는 회로/펌웨어 문제가 아닌
        모터 물성으로, 학습 환경에서는 \texttt{max\_speed}를 30\,\%
        이상에서 운영할 것을 권장한다.
\end{itemize}

\section{2026-05-20 --- Web Bluetooth와 HTTPS 서버 사용 사유}

\subsection{관찰}

\texttt{Web/ks01.html}은 \texttt{file://} 프로토콜로 직접 열어도 BLE
연결이 가능했다. 반면 정식 학습용 \texttt{Web/main.html}은
\texttt{file://}에서 연결되지 않고, \texttt{http://localhost} 또는
\texttt{https://} 호스팅 환경에서만 동작한다.

\subsection{원인 --- Web Bluetooth가 아니라 ES module의 same-origin 정책}

직관적으로는 ``Web Bluetooth가 \texttt{file://}에서 secure context로
인정되지 않기 때문''으로 생각하기 쉽지만, 실제 원인은 다르다. 두
파일의 스크립트 로딩 방식이 결정적으로 다르다.

\begin{table}[h]
\centering
\small
\begin{tabularx}{\linewidth}{@{}lXX@{}}
\toprule
파일 & 스크립트 로딩 방식 & \texttt{file://}에서 동작? \\
\midrule
\texttt{ks01.html} & 모든 코드가 인라인 \texttt{<script>...</script>} & 정상 동작 \\
\texttt{main.html} & \texttt{<script type="module" src="main.js">}, \texttt{main.js}는 \texttt{bluetooth.js} 등 다수 모듈을 \texttt{import} & 차단 (모듈 fetch 실패) \\
\bottomrule
\end{tabularx}
\caption{스크립트 로딩 방식과 \texttt{file://} 동작 가부.}
\end{table}

Chromium 기반 브라우저는 \texttt{file://}을 secure context로 인정하므로
\texttt{navigator.bluetooth} API 자체는 \texttt{file://}에서도 호출
가능하다. 그러나 \texttt{<script type="module">}은 \emph{같은 출처
(same-origin) 정책}을 따르는데, 모든 \texttt{file://} URL은 같은
디렉터리에 있더라도 \texttt{origin: null}로 취급되어 서로 다른 출처로
간주된다. 따라서 \texttt{main.js}를 비롯한 ES 모듈 fetch가 CORS
정책으로 차단되고, 콘솔에는 다음과 같은 오류가 기록된다.

\begin{minted}{text}
Access to script at 'file:///.../main.js' from origin 'null'
has been blocked by CORS policy: Cross origin requests are
only supported for protocol schemes: http, data, isolated-app,
chrome-extension, chrome-untrusted, https, chrome.
\end{minted}

ES 모듈이 로드되지 않으면 \texttt{BluetoothManager} 코드는 실행될
기회조차 없다. 즉 사용자가 본 ``BLE가 연결되지 않는'' 증상은 Web
Bluetooth API의 보안 제한이 아니라 \emph{BLE 코드가 아예 실행되지
않은} 결과이다. 증상이 같아 보였을 뿐이다.

\begin{caution}
\texttt{ks01.html}이 \texttt{file://}에서 동작하는 것은 데스크톱 Chrome
환경에 한정될 가능성이 크다. Android Chrome은 \texttt{file://}에서 Web
Bluetooth를 더 엄격히 제한하는 경우가 있으며, 모바일 배포를 가정한다면
이 경로에 의존해서는 안 된다.
\end{caution}

\subsection{해결 옵션과 권장}

\begin{table}[h]
\centering
\small
\begin{tabularx}{\linewidth}{@{}lXl@{}}
\toprule
방법 & 설명 & 권장 \\
\midrule
A. \texttt{python -m http.server} & \texttt{localhost}는 secure context로 인정되고 ES 모듈도 같은 출처에서 fetch 가능 & 개발 단계 \\
B. HTTPS 호스팅 (GitHub Pages 등) & 어디서나 동작, 모바일 포함. \texttt{file://} 의존 사라짐 & 학생 배포 \\
C. 단일 HTML로 번들링(esbuild/rollup) & \texttt{main.js}와 모듈을 인라인 IIFE로 합쳐 \texttt{ks01.html}과 같은 형태 & 오프라인 단일 파일이 꼭 필요한 경우 \\
D. ES module 제거, 전역 객체로 통합 & \texttt{import/export} 제거 후 인라인 \texttt{<script>}로 합치기 & 권장하지 않음 (모듈성 손실 큼) \\
\bottomrule
\end{tabularx}
\caption{\texttt{main.html}의 \texttt{file://} 미동작에 대한 해결 옵션.}
\end{table}

\begin{keypoint}
학생 배포를 전제하면 \textbf{B(HTTPS 호스팅)}\,가 가장 깔끔하다. 모바일
브라우저 호환성, 즐겨찾기/PWA 설치, 자동 업데이트 측면 모두 유리하다.
\texttt{ks01.html}처럼 단일 파일 배포가 정책적으로 요구되는 경우에는
\textbf{C(번들링)}\,가 정공법이다. 어느 쪽이든 \texttt{file://} 직접
열기에 의존하는 워크플로우는 장기적으로 권장되지 않는다.
\end{keypoint}

\section{잔여 위험 및 후속 작업}

\subsection{Web 측 \texttt{saveSystemSettings()} 응답 검증}

SYS\_SET 송신 후 무조건 성공 alert를 띄우는 현행 동작은 펌웨어에서
실제로 저장 실패한 경우를 사용자에게 가린다. 응답 대기 명령으로
승격하여 \texttt{"1"}/\texttt{"0"} 또는 \texttt{"OK"}/\texttt{"ERROR"}
구분 후 분기 alert를 띄우는 작업이 권장된다.

\subsection{UART baud rate}

가장 본질적인 처리량 천장은 BLE 모듈 $\leftrightarrow$ Pico 사이의 9600
baud UART이다. HM-10/BT05의 AT 명령(\texttt{AT+BAUD8} 등)으로 38400
또는 115200 baud로 상향하고 \texttt{Pico/main.py}의
\texttt{UART\_BAUDRATE}를 동기화하면 byte당 전송 시간을 4--12배 단축할
수 있다. BLE 모듈 펌웨어 설정과 Pico 코드 동시 수정이 필요하므로 별도
작업으로 분리한다.

\subsection{호스팅 환경 정비}

\texttt{Web/main.html}의 ES module 의존을 그대로 두면서 학생 배포를
원활하게 하려면 GitHub Pages 등 HTTPS 호스팅을 정식 채택할 필요가 있다.
도메인/배포 자동화/PWA manifest 추가가 후속 작업 후보이다.

\subsection{권장 회귀 테스트}

\begin{enumerate}
  \item \textbf{카운트다운} --- \texttt{LED\_OFF} 5개 연속 소등. 명령
        사이 체감 지연이 거의 사라져야 한다.
  \item \textbf{LED 패턴 + 거리 센서 혼합} --- \texttt{[v0 v1 v2 v3 v4]}
        직후 \texttt{DISTANCE}. 거리값이 \texttt{1} 같은 LED 응답으로
        오염되지 않는지 확인.
  \item \textbf{부저 사이렌} --- \texttt{BUZZER\_ON,400,0.5} 와
        \texttt{BUZZER\_ON,800,0.5}의 반복. 타이밍 정확성이 유지되어야
        한다.
  \item \textbf{시간 제한 모터} --- \texttt{DC\_tFORWARD,N},
        \texttt{DC\_tBACKWARD,N}의 N초 동기화가 유지되는지 확인.
  \item \textbf{DC 모터 양방향 대칭} --- \texttt{max\_speed}를 30, 50,
        80, 100\,\%로 바꿔가며 전진/후진 속도가 거의 같게 느껴지는지
        확인.
  \item \textbf{대시보드 \texttt{max\_speed} 영구 저장} --- 슬라이더
        조정 $\rightarrow$ 업로드 $\rightarrow$ 다른 탭 이동 $\rightarrow$ 대시보드 복귀 시
        값이 유지되는지 확인. Pico 재부팅 후에도 유지되어야 한다.
\end{enumerate}

\section{변경 파일 요약}

\begin{table}[h]
\centering
\small
\begin{tabularx}{\linewidth}{@{}llX@{}}
\toprule
일자 & 파일 & 변경 요지 \\
\midrule
05-19 & \texttt{Web/commandexecutor.js} & \texttt{FIRE\_AND\_FORGET\_HEADS} 집합 도입, \texttt{sendCommand()}에서 분기 처리, cooldown 분리(40 / 20\,ms). \\
05-19 & \texttt{Pico/main.py} & \texttt{NO\_RESPONSE\_CMDS} 클래스 변수, \texttt{\_read\_uart\_line()} 폴링 방식으로 교체, \texttt{\_needs\_response()} 도입, \texttt{run()}에서 조건부 \texttt{\_send\_response()}. \\
\midrule
05-20 & \texttt{Pico/main.py} & \texttt{boot()}의 OLED 호출 None 가드 추가. 05-19 BLE 최적화 재적용. \\
05-20 & \texttt{Pico/process\_data.py} & 누락되어 있던 \texttt{\_handle\_sys\_set()} 메서드 헤더 복구. \\
05-20 & \texttt{Pico/dcmotor.py} & \texttt{DIR+PWM} 단일 변조에서 양 핀 모두 PWM으로 전환. \texttt{dc\_forward}/\texttt{dc\_backward}가 한쪽 핀을 \texttt{PWM\_MAX} 고정, 다른 쪽을 \texttt{PWM\_MAX-speed}로 변조하여 양방향 brake-modulated 대칭 출력. \\
\bottomrule
\end{tabularx}
\caption{2026-05-19 \textendash{} 05-20 변경 범위.}
\end{table}

\begin{table}[h]
\centering
\small
\begin{tabularx}{\linewidth}{@{}llX@{}}
\toprule
일자 & 커밋 & 메시지 \\
\midrule
05-19 & \texttt{7e5b93f} & Optimize BLE pipeline: fire-and-forget for output commands \\
05-19 & \texttt{110bb27} & Revert ``Optimize BLE pipeline: fire-and-forget for output commands'' \\
05-20 & \texttt{600243a} & Re-apply BLE pipeline optimization with OLED boot crash fix \\
05-20 & \texttt{de9cffc} & Fix DC motor: forward PWM asymmetry and missing SYS\_SET handler \\
05-20 & \texttt{5780485} & Drive H-bridge with dual PWM for symmetric DC motor output \\
\bottomrule
\end{tabularx}
\caption{관련 커밋 (\texttt{main} 브랜치).}
\end{table}

\end{document}
